\chapter{Bleeding Control}\label{ch:bleeding}

\section{Understanding Bleeding Types}

Bleeding (hemorrhage) can be life-threatening if not controlled promptly. Understanding the type of bleeding helps determine appropriate management.

\subsection{Capillary Bleeding}

Slow, oozing blood from small vessels in superficial wounds. Usually minor and easy to control with direct pressure and a bandage.

\subsection{Venous Bleeding}

Dark red or maroon blood that flows steadily from the wound. Caused by damage to veins (blood returning to the heart). Generally less dangerous than arterial bleeding but requires prompt attention.

\subsection{Arterial Bleeding}

Bright red blood that pulses or spurts from the wound, often forcefully timed with each heartbeat. Caused by damage to arteries (blood leaving the heart). This is the most dangerous type and can lead to fatal blood loss within minutes.

\textbf{WARNING:} Any uncontrolled bleeding should be treated as potentially life-threatening. Seek emergency medical help immediately for significant bleeding.

\section{Assessing Blood Loss}

Rapid assessment of blood loss severity guides treatment urgency.

\begin{center}
\small
\begin{tabularx}{\textwidth}{@{}>{\raggedright\arraybackslash}l>{\raggedright\arraybackslash}X>{\raggedright\arraybackslash}X>{\raggedright\arraybackslash}X@{}}
\hline
\textbf{Class} & \textbf{Blood Loss (\%)} & \textbf{Volume (mL)} & \textbf{Signs \& Symptoms} \\
\hline
I & Up to 15\% & <750 mL & Minimal changes; may be asymptomatic \\
II & 15-30\% & 750-1500 mL & Tachycardia (>100 bpm), tachypnea, anxiety, delayed capillary refill \\
III & 30-40\% & 1500-2000 mL & Marked tachycardia (>120 bpm), hypotension, altered mental status, rapid shallow breathing \\
IV & >40\% & >2000 mL & Profound shock: minimal or no urine output, significant confusion/unconsciousness, weak/no palpable pulse \\
\hline
\end{tabularx}
\end{center}

\quicknote{\textbf{Normal Adult Blood Volume}: Approximately 5-6 liters (about 70 mL/kg body weight).}

\section{Basic Hemostasis Techniques}

The primary method for controlling external bleeding is direct pressure.

\subsection{Direct Pressure}

\begin{enumerate}[leftmargin=*]
\item Put on disposable gloves if available
\item Apply a clean dressing, sterile gauze, or cloth directly onto the wound
\item Press firmly and continuously on the wound for at least 5 minutes without checking -- lifting to check disrupts clot formation
\item If blood soaks through, DO NOT remove the soaked dressing; add more layers on top and continue pressure
\item Once bleeding slows, secure the dressing with a bandage
\end{enumerate}
\subsection{Elevation}

Raise the injured limb above the level of the heart (if no fracture suspected) to reduce blood flow to the area. Combine with direct pressure for maximum effect.

\subsection{Pressure Points}

If direct pressure and elevation are insufficient, apply pressure to major arteries against underlying bone.

\begin{itemize}[leftmargin=*]
  \item \textbf{Brachial artery}: Inner upper arm (for arm injuries)
  \item \textbf{Femoral artery}: Groin/thigh crease (for leg injuries)
  \item \textbf{Radial/Wrist}: Not typically effective for major bleeding control
\end{itemize}

\textbf{WARNING:} \textbf{Pressure points are temporary measures only.} They do not replace direct pressure at the wound site. Use until professional help arrives.

\section{Tourniquet Application}

A tourniquet is a last resort for life-threatening extremity bleeding that cannot be controlled by other methods \cite{aha_arc_2024}.

\subsection{When to Apply a Tourniquet}

\begin{itemize}[leftmargin=*]
  \item Severe, spurting arterial bleeding from an arm or leg
  \item Direct pressure has failed after several attempts
  \item Amputation injury with massive bleeding
  \item Multiple casualties where you must move to the next victim
\end{itemize}

\textbf{WARNING:} Tourniquets should NEVER be used for facial, neck, torso, or groin injuries. They should also not be used on joints (knee/elbow).

\begin{enumerate}[leftmargin=*]
\item Apply 2-3 inches above the wound (closer to the heart), NOT over a joint
\item Avoid embedded objects if possible (place tourniquet above them)
\item Wrap tightly enough to stop bleeding completely
\item Note the time of application clearly (on forehead or on tourniquet)
\end{enumerate}
\textbf{WARNING:} \textbf{Important: Once applied, a tourniquet should NOT be loosened or removed until medical professionals take over. Loosening it can restart severe bleeding. Always document the exact time of application for medical responders.}

\section{Managing Embedded Objects}\label{sec:embedded_objects}

Do not remove large objects impaled in wounds (sticks, glass, metal, etc.), as they may be plugging the bleeding vessel and removing them could cause catastrophic hemorrhage.

\begin{enumerate}[leftmargin=*]
\item Control bleeding around the object with dressings placed adjacent to it (not on top)
\item Stabilize the object with bulky dressings to prevent movement
\item Seek immediate medical assistance
\end{enumerate}
\textbf{WARNING:} \textbf{NEVER} attempt to remove an embedded object. Leave removal to trained medical personnel.

\section{Nosebleeds (Epistaxis)}

Nosebleeds are common and usually benign but can be alarming.

\begin{enumerate}[leftmargin=*]
\item Have the person sit upright and lean slightly forward (NOT backward)
\item Pinch the soft part of the nose between thumb and forefinger
\item Maintain firm pressure for at least 10 minutes continuously without releasing
\item After 10 minutes, release gently and check if bleeding has stopped; repeat if necessary
\item After bleeding stops, advise blowing the nose gently or avoiding strenuous activity for 24 hours
\end{enumerate}
\textbf{WARNING:} \textbf{DO NOT tilt head backward, as this allows blood to run down the throat, which can cause nausea or vomiting. Persistent nosebleeds (more than 20 minutes despite pressure) require medical evaluation.}

\section{Common Mistakes to Avoid}

\begin{itemize}[leftmargin=*]
  \item Removing soaked dressings -- adds more layers instead
  \item Lifting to check before 5 continuous minutes of pressure
  \item Using dirty materials that could introduce infection
  \item Applying tourniquets incorrectly (on joints, too loose, not timed)
  \item Rushing medical transport when bleeding is controlled
\end{itemize}

\section{When to Call Emergency Services}

Call emergency services (112 or 108 in India; check your local number elsewhere) if:

\begin{itemize}[leftmargin=*]
  \item Bleeding does not stop after 10-15 minutes of direct pressure
  \item Blood is spurting brightly red (arterial bleeding)
  \item A large object is embedded in the wound
  \item There is significant head, neck, chest, or abdominal trauma
  \item The person shows signs of shock (pale skin, rapid pulse, confusion, weakness)
  \item The wound was caused by a serious accident (vehicle collision, fall from height)
  \item There is animal bite or puncture wound requiring tetanus/antibiotics
\end{itemize}

\medskip\noindent\textbf{Hemorrhage Control Summary.}
\begin{enumerate}
  \item Apply direct pressure on the wound with clean material
  \item Raise the injured area above heart level if possible
  \item Add more layers if blood soaks through
  \item Do not remove embedded objects
  \item Use tourniquet ONLY for life-threatening extremity bleeding when other methods fail
  \item Note time of tourniquet application
  \item Seek emergency medical help for persistent or severe bleeding
\end{enumerate}

\section{Review Questions}

\begin{enumerate}[leftmargin=*]
\item How is arterial bleeding distinguished from venous bleeding, and why does it matter?
\item Why should a blood-soaked dressing never be removed?
\item When is a tourniquet appropriate, and what must you always do when you apply one?
\item A large object is embedded in a wound. What should you do and why?
\item What is the correct position for a person with a nosebleed, and how long should pressure be held?
\end{enumerate}
\section{Skill Practice: Bleeding Control}

\begin{enumerate}[leftmargin=*]
\item Don gloves; inspect the simulated wound while maintaining scene safety
\item Apply a dressing and firm direct pressure for a timed 5 minutes without peeking
\item Apply additional layers over a soaked dressing
\item Demonstrate limb elevation with a simulated fracture-free injury
\item Apply a training tourniquet 2--3 inches above a simulated amputation; tighten until a distal pulse disappears; note the time
\item Practice stabilizing an embedded object with bulky dressings on both sides
\item Report the situation and interventions as if handing over to responders
\end{enumerate}
\index{bleeding}
\index{hemorrhage}
\index{direct pressure}
\index{tourniquet}
\index{embedded objects}
\index{epistaxis}
\index{shock}
\index{pressure points}
\index{venous bleeding}
\index{arterial bleeding}
\index{capillary bleeding}